\begin{frame}{자주 하는 실수: 리뷰가 감상평에 머무는 것}
  \begin{itemize}
    \item ``발표 잘 들었습니다'', ``결과가 인상적이네요'' ---
      체크리스트를 \textbf{형식적으로 채웠을 뿐}, 발표팀에게
      실제로 도움이 되는 정보를 전혀 주지 않는다.
    \item \textbf{구체적인 주차·개념을 근거로 든 질문}만이
      ``반박 가능''하다 --- 발표팀이 답하려면 실제로 자신의
      결과를 다시 들여다봐야 하기 때문.
  \end{itemize}
  \vskip0.6em
  \begin{block}{채점의 축}
    peer-review 점수가 체크리스트를 ``채웠는가''가 아니라
    \textbf{``그 안의 질문이 얼마나 구체적인가''}로 매겨지는
    이유가 여기에 있다.
  \end{block}
\end{frame}
